\begin{conclusion}[指数对数坐标下的 Lee--Yang 微分闭合：\texorpdfstring{$\lambda(s)=\lambda_+(e^s)$}{lambda(s)} 的一阶代数方程与椭圆核一致性]\label{con:xi-terminal-zm-leyang-exponential-time-algebraic-differential-closure}
沿用谱四次
\[
\Pi(\lambda,y):=\lambda^4-\lambda^3-(2y+1)\lambda^2+\lambda+y(y+1)=0
\]
及其在 \(y=1\) 邻域满足 \(\lambda_+(1)=2\) 的解析特征支 \(\lambda_+(y)\)。
在 \(s=0\) 的邻域定义对数参数化
\[
\lambda(s):=\lambda_+(e^s),
\qquad
\dot\lambda:=\frac{\dd \lambda}{\dd s}.
\]
则 \(\lambda(s)\) 满足显式一阶代数微分闭式
\begin{equation}\label{eq:xi-terminal-zm-leyang-exponential-time-lambda-ode}
\boxed{\;
\begin{aligned}
0={}&(16\lambda^4+7\lambda^3-9\lambda^2+\lambda+1)\,\dot\lambda^{2}
(-16\lambda^5-4\lambda^4+20\lambda^3-5\lambda+1)\,\dot\lambda \\
&\qquad\qquad+(4\lambda^6-8\lambda^4+\lambda^3+4\lambda^2-\lambda).
\end{aligned}\;}
\end{equation}
并且将 \eqref{eq:xi-terminal-zm-leyang-exponential-time-lambda-ode} 视为关于 \(\dot\lambda\) 的二次方程，其判别式满足严格因式分解
\begin{equation}\label{eq:xi-terminal-zm-leyang-exponential-time-lambda-ode-discriminant}
\boxed{\;
\Delta_{\dot\lambda}(\lambda)
=\bigl(4\lambda^3-4\lambda+1\bigr)\bigl(2\lambda^3+2\lambda^2-\lambda+1\bigr)^2\; }.
\end{equation}
因此在函数域层面成立
\begin{equation}\label{eq:xi-terminal-zm-leyang-exponential-time-function-field}
\QQ(\lambda,\dot\lambda)=\QQ\!\left(\lambda,\sqrt{4\lambda^3-4\lambda+1}\right),
\end{equation}
从而指数对数坐标的微分闭合与椭圆核
\begin{equation}\label{eq:xi-terminal-zm-leyang-exponential-time-elliptic-kernel}
\mathcal E:\quad w^2=4\lambda^3-4\lambda+1
\end{equation}
处于同一函数域（与结论 \ref{con:xi-terminal-zm-elliptic-normalization} 的椭圆化一致）。
更进一步，外参 \(y=e^s\) 可由 \((\lambda,\dot\lambda)\) 有理重构为
\begin{equation}\label{eq:xi-terminal-zm-leyang-exponential-time-y-reconstruction}
\boxed{\;
y=
\frac{\dot\lambda\,(4\lambda^3-3\lambda^2-2\lambda+1)-2\lambda(\lambda-1)^2(\lambda+1)}{-2\lambda^2+4\lambda\dot\lambda+1}\; }.
\end{equation}
\end{conclusion}
\begin{proof}[证明要点（隐式微分与消元）]
对恒等式 \(\Pi(\lambda(s),e^s)\equiv0\) 微分得
\(
\partial_\lambda\Pi(\lambda,y)\,\dot\lambda+\partial_y\Pi(\lambda,y)\,\dot y=0
\)
且 \(\dot y=y\)。
由 \(\Pi(\lambda,y)=0\) 与
\(
\partial_\lambda\Pi(\lambda,y)\,\dot\lambda+y\,\partial_y\Pi(\lambda,y)=0
\)
得到关于 \(y\) 的两条二次代数约束；消去 \(y^2\) 得到 \(y\) 的有理表达式 \eqref{eq:xi-terminal-zm-leyang-exponential-time-y-reconstruction}，
代回即得 \eqref{eq:xi-terminal-zm-leyang-exponential-time-lambda-ode}。
判别式分解 \eqref{eq:xi-terminal-zm-leyang-exponential-time-lambda-ode-discriminant} 为直接计算并在 \(\ZZ[\lambda]\) 中因式分解。
函数域恒等式 \eqref{eq:xi-terminal-zm-leyang-exponential-time-function-field} 由判别式的唯一非平凡平方自由部 \((4\lambda^3-4\lambda+1)\) 立即推出；
而 \eqref{eq:xi-terminal-zm-leyang-exponential-time-elliptic-kernel} 与结论 \ref{con:xi-terminal-zm-elliptic-normalization} 的椭圆模型一致。
可复现核验脚本：\texttt{scripts/exp\_fold\_zm\_leyang\_exponential\_time\_differential\_closure\_audit.py}。证毕。
\end{proof}

\begin{conclusion}[\texorpdfstring{$\psi(s)=\log\lambda_+(e^s)-\log2$}{psi(s)} 的微分代数性与非 \texorpdfstring{$D$}{D}-有限性]\label{con:xi-terminal-zm-psi-differential-algebraic-non-dfinite}
令
\[
\psi(s):=\log\lambda_+(e^s)-\log2.
\]
则 \(\psi\) 在 \(s=0\) 的解析芽满足：
\begin{enumerate}
  \item \(\psi\) 为微分代数函数：其满足某个系数属 \(\QQ(s)\) 的非平凡代数微分方程；
  \item \(\psi\) 非 \(D\)-有限：不存在首项系数为多项式的非零线性常微分算子湮灭 \(\psi\)。
  因而泰勒系数列 \(\bigl(\psi^{(n)}(0)\bigr)_{n\ge 0}\) 非 \(P\)-递归（非 holonomic）。
\end{enumerate}
\end{conclusion}
\begin{proof}[证明要点]
第一条由结论 \ref{con:xi-terminal-zm-leyang-exponential-time-algebraic-differential-closure} 给出的 \(\lambda(s)=\lambda_+(e^s)\) 的一阶代数微分闭式推出：
\(\lambda\) 因而微分代数；又 \(\psi'=\dot\lambda/\lambda\) 为 \((\lambda,\dot\lambda)\) 的有理函数且 \(\lambda(0)=2\neq0\)，故 \(\psi\) 的解析芽同属微分代数闭包。
\par
第二条利用奇点集合的无限性。
由结论 \ref{con:xi-terminal-zm-discriminant-leyang}，\(\lambda_+(y)\) 的有限分歧值集合包含某个满足 \(y_\star\neq0\) 的代数分歧值（例如 Lee--Yang 三次 \(P_{\mathrm{LY}}(y)=0\) 的任一根）。
取任一对数支 \(s_\star:=\log y_\star\)，则指数映射的周期性给出无限前像族
\[
s_\star+2\pi i\,\ZZ.
\]
在每个点 \(s=s_\star+2\pi i k\) 处，复合函数 \(s\mapsto \lambda_+(e^s)\) 继承 \(y=y_\star\) 的代数分歧而出现不可去奇性，故 \(\psi\) 在 \(\CC\) 上具有无穷多个奇点。
另一方面，若 \(\psi\) 为 \(D\)-有限，则其满足某个首项系数为多项式的线性常微分方程，故其解析延拓的可能奇点仅能位于该首项系数的零点集合（有限集）与无穷远点。
两者矛盾，故 \(\psi\) 非 \(D\)-有限。
最后，\(D\)-有限解析芽的泰勒系数列必为 \(P\)-递归为标准等价，因而 \(\bigl(\psi^{(n)}(0)\bigr)\) 非 \(P\)-递归。证毕。
\end{proof}

\begin{conclusion}[主特征支局部喷射的 \(3\)-局部整性：\texorpdfstring{$\lambda_+(1+z)\in\ZZ[1/3][[z]]$}{lambda+(1+z) in Z[1/3]} 与 \texorpdfstring{$2^n\psi^{(n)}(0)\in\ZZ[1/3]$}{2^n psi^(n)(0) in Z[1/3]}]\label{con:xi-terminal-zm-lambda-plus-3local-denominators}
记
\[
\lambda_+(1+z)=2+\sum_{n\ge 1}a_n z^n\in\QQ[[z]],
\qquad
\psi(s)=\sum_{n\ge 1}\frac{\psi^{(n)}(0)}{n!}\,s^n.
\]
则：
\begin{enumerate}
  \item 对一切 \(n\ge 1\) 有 \(a_n\in\ZZ[1/3]\)，亦即 \(\lambda_+(1+z)\in\ZZ[1/3][[z]]\)；
  \item 对一切 \(n\ge 1\) 有 \(2^n\psi^{(n)}(0)\in\ZZ[1/3]\)，从而 \(\psi^{(n)}(0)\in\ZZ[1/6]\)。
\end{enumerate}
\end{conclusion}
\begin{proof}[证明要点（\(p\)-进隐函数与导数闭合）]
由 \(\Pi(2,1)=0\) 且
\(
\partial_\lambda\Pi(2,1)=9
\)，
对任一素数 \(p\neq3\)，\(\partial_\lambda\Pi(2,1)\) 在 \(\ZZ_p\) 中可逆。
由 \(p\)-进隐函数定理，存在唯一 \(\lambda_p(z)\in\ZZ_p[[z]]\) 使
\(
\Pi(\lambda_p(z),1+z)\equiv0
\)
且 \(\lambda_p(0)=2\)。
而 \(\lambda_+(1+z)\in\QQ[[z]]\) 亦满足同一初值条件，故在 \(\QQ_p[[z]]\) 中与 \(\lambda_p\) 一致；
于是 \(a_n\in\ZZ_p\) 对一切 \(p\neq3\) 成立，从而 \(a_n\in\ZZ[1/3]\)。
\par
再记 \(\lambda(s)=\lambda_+(e^s)\)。
由 \(e^s\) 在 \(s=0\) 处各阶导数皆为 \(1\)，复合求导（Fa\`a di Bruno）表明 \(\lambda^{(n)}(0)\) 为 \(\{\lambda_+^{(k)}(1)\}_{k\le n}\) 的整系数组合；
而 \(\lambda_+^{(k)}(1)=k!\,a_k\in\ZZ[1/3]\)，故 \(\lambda^{(n)}(0)\in\ZZ[1/3]\)。
又 \(\psi=\log(\lambda/2)\) 且 \(\lambda(0)=2\)，于是 \(\psi^{(n)}(0)\) 为 \(\{\lambda^{(k)}(0)/2\}_{k\le n}\) 的整系数多项式（来自对数导数的递推），
每项分母至多贡献 \(2^n\)，从而 \(2^n\psi^{(n)}(0)\in\ZZ[1/3]\)。证毕。
\end{proof}

\begin{conclusion}[有限阶局部喷射唯一确定谱四次：\(13\)-阶芽锁定 \texorpdfstring{$\Pi(\lambda,y)$}{Pi(lambda,y)} 与 Lee--Yang 三次]\label{con:xi-terminal-zm-local-jet-rigidly-determines-pi}
在二元多项式集合中考虑逆问题：设
\(
Q(\lambda,y)\in\QQ[\lambda,y]
\) 满足 \(\deg_\lambda Q\le4\)、\(\deg_y Q\le2\)，并且 \(\lambda^4\) 的系数归一为 \(1\)。
若已知主特征支 \(\lambda_+(y)\) 在 \(y=1\) 的 \(13\)-阶喷射（即 \(a_1,\dots,a_{13}\) 连同 \(\lambda_+(1)=2\)），
则满足
\[
Q(\lambda_+(y),y)\equiv 0 \qquad\text{于}\qquad \QQ[[y-1]]/(y-1)^{14}
\]
的上述 \(Q\) 唯一存在，且必为谱四次 \(\Pi\) 本身。
因此由恒等式
\(
\mathrm{Disc}_{\lambda}(\Pi)=-y(y-1)P_{\mathrm{LY}}(y)
\)
所定义的 Lee--Yang 三次 \(P_{\mathrm{LY}}(y)\) 亦被该有限喷射唯一决定。
\end{conclusion}
\begin{proof}[证明要点（线性化与可逆性）]
将一般 \(Q\) 写为
\[
Q(\lambda,y)=\lambda^4+\sum_{(i,j)\neq(4,0)}c_{ij}\lambda^i y^j,
\qquad 0\le i\le4,\ 0\le j\le2,
\]
未知量为 \(14\) 个系数 \(c_{ij}\in\QQ\)。
代入 \(y=1+z\) 与
\(
\lambda=\lambda_+(1+z)=2+\sum_{n=1}^{13}a_n z^n+O(z^{14})
\)，
并取模 \(z^{14}\)，条件 \(Q(\lambda_+(1+z),1+z)\equiv0\pmod{z^{14}}\) 等价于 \(z^0,\dots,z^{13}\) 的 \(14\) 个系数同时为零，
从而得到一个 \(14\times14\) 的线性方程组。
直接计算可得其系数矩阵行列式非零，故解唯一存在。
另一方面，\(\Pi\) 显然给出一个满足条件的解，故唯一解必为 \(\Pi\)。
由判别式恒等式 \(\mathrm{Disc}_{\lambda}(\Pi)=-y(y-1)P_{\mathrm{LY}}(y)\) 立得 \(P_{\mathrm{LY}}\) 的唯一性。证毕。
可复现核验脚本：\texttt{scripts/exp\_fold\_zm\_leyang\_exponential\_time\_differential\_closure\_audit.py}。
\end{proof}

